
Une fois le backlog du produit finalisé, nous avons organisé la réunion de planification de sprint. L'objectif de cette réunion était de construire le backlog de sprint en se basant sur le backlog de produit défini par le Product Owner. À l'issue de cette réunion, nous avons identifié les durées prévisionnelles du travail à effectuer pour chaque sprint. Pour ce projet, nous avons divisé le travail en trois sprints, chacun d'une durée de 3 semaines, regroupés en une seule release. La figure 2.3 illustre la répartition des sprints pour notre système.

\medskip

\textbf{Plan des sprints :}

\begin{itemize}
    \item \textbf{Sprint 1 : Collecte de données basée sur caméra}
    \begin{itemize}
        \item \textbf{Objectif} : Mettre en place un système de collecte de données à partir des moniteurs des signes vitaux en utilisant des caméras et des techniques de reconnaissance optique de caractères (OCR).
        \item \textbf{Deliverables} :
            \begin{itemize}
                \item Intégration de la caméra IP avec le système.
                \item Pipeline OCR pour extraire les valeurs vitales (fréquence cardiaque, pression artérielle, SpO2).
            \end{itemize}
        \item \textbf{Durée} : 3 semaines.
        \item \textbf{Tâches clés} :
            \begin{itemize}
                \item Configuration des caméras IP (montage fixe, résolution 1080p).
                \item Installation et configuration des logiciels OCR (OpenCV, Tesseract).
                \item Développement du pipeline de prétraitement et extraction d'images.
            \end{itemize}
        \item \textbf{Critères d'acceptation} :
            \begin{itemize}
                \item Le pipeline OCR extrait les valeurs vitales avec une précision d'au moins 95\%, validée par comparaison avec 10 lectures manuelles.
                \item Les images sont traitées avec une latence inférieure à 1 seconde par frame dans 100\% des tests.
                \item Le pipeline est compatible avec les moniteurs Mindray et Edan, testé sur au moins 5 scénarios différents.
            \end{itemize}
        \item \textbf{Risques} :
            \begin{itemize}
                \item Mauvaise qualité d'image affectant la précision de l'OCR.
                \item Solution : Utilisation de caméras haute résolution et optimisation du prétraitement des images avec OpenCV.
            \end{itemize}
    \end{itemize}

    \item \textbf{Sprint 2 : Validation des données et gestion des alertes}
    \begin{itemize}
        \item \textbf{Objectif} : Valider les données collectées et développer un système de gestion des alertes basé sur des seuils prédéfinis.
        \item \textbf{Deliverables} :
            \begin{itemize}
                \item Validation des données OCR par rapport aux données réelles.
                \item Module de gestion des alertes avec seuils prédéfinis.
                \item Notifications push via Firebase Cloud Messaging (FCM).
            \end{itemize}
        \item \textbf{Durée} : 3 semaines.
        \item \textbf{Tâches clés} :
            \begin{itemize}
                \item Tests de validation des données extraites (comparaison avec lectures manuelles).
                \item Développement du module de détection des anomalies basé sur des seuils.
                \item Intégration de FCM pour les notifications push sur mobile.
            \end{itemize}
        \item \textbf{Critères d'acceptation} :
            \begin{itemize}
                \item Les données extraites correspondent aux valeurs réelles avec une précision de 98\% dans 20 tests.
                \item Les alertes sont déclenchées en moins de 2 secondes lorsque les seuils sont dépassés (e.g., fréquence cardiaque > 120 bpm).
                \item Les notifications push sont reçues avec un taux de succès de 99\% sur Android et iOS dans un environnement de test.
            \end{itemize}
        \item \textbf{Dépendances} : Données collectées et pipeline OCR fonctionnel du Sprint 1.
        \item \textbf{Risques} :
            \begin{itemize}
                \item Retards dans la validation des données en raison d'erreurs OCR.
                \item Solution : Ajout de tests unitaires et manuels supplémentaires pour l'OCR.
            \end{itemize}
    \end{itemize}

    \item \textbf{Sprint 3 : Interface utilisateur et stockage des données}
    \begin{itemize}
        \item \textbf{Objectif} : Développer une interface utilisateur pour visualiser les données et alertes, et mettre en place le stockage des données.
        \item \textbf{Deliverables} :
            \begin{itemize}
                \item Tableau de bord web (React) pour afficher les données vitales et alertes.
                \item Application mobile (Flutter) pour les notifications et la consultation.
                \item Base de données MySQL pour stocker les données et historiques.
            \end{itemize}
        \item \textbf{Durée} : 3 semaines.
        \item \textbf{Tâches clés} :
            \begin{itemize}
                \item Conception et développement du tableau de bord web avec React.
                \item Développement de l'application mobile Flutter avec intégration FCM.
                \item Configuration de la base de données MySQL pour le stockage des données et alertes.
            \end{itemize}
        \item \textbf{Critères d'acceptation} :
            \begin{itemize}
                \item Le tableau de bord affiche les données vitales et alertes en temps réel avec un rafraîchissement toutes les 5 secondes, validé par 5 utilisateurs médicaux.
                \item L'application mobile permet la consultation des alertes avec une ergonomie validée par 3 infirmiers (aucun problème d'utilisabilité signalé).
                \item Les données sont stockées dans MySQL avec 100\% d'intégrité (aucune perte de données dans 50 tests).
            \end{itemize}
        \item \textbf{Dépendances} : Données validées et système d'alertes du Sprint 2.
        \item \textbf{Risques} :
            \begin{itemize}
                \item Complexité dans l'intégration des interfaces web et mobile.
                \item Solution : Prototypage préliminaire et tests d'utilisabilité itératifs.
            \end{itemize}
    \end{itemize}
\end{itemize}